%%%%%%%%%%  scan idx 450  |  folio 437  %%%%%%%%%%
On sait (EGA II 8.1.8) que \(Y'\) est un sous-schéma fermé de \(X'\) défini par un idéal de \(\cO_X\), canoniquement isomorphe à \(\cO_{X'}(1)\); l'immersion \(j:Y'\to X\) étant définie par un idéal inversible, est une immersion fermée régulière de codimension \(1\). On note \(J\) l'idéal de \(\cO_X\) qui définit \(Y\), et \(J'\) l'idéal de \(\cO_{X'}\) qui définit \(Y'\); on a donc \(J'\simeq J.\cO_{X'}\simeq\cO_{X'}(1)\). Enfin, on note \(f:X'\to X\), et \(g:Y'\to Y\) les morphismes structuraux. On a donc le diagramme commutatif:
\begin{equation}\tag{3.1.1}
\begin{tikzcd}[column sep=large,row sep=large]
Y' \arrow[r,"j"] \arrow[d,"g"'] & X' \arrow[d,"f"]\\
Y \arrow[r,"i"'] & X.
\end{tikzcd}
\end{equation}

Par définition, \(X'=\Proj(S)\), où \(S=\bigoplus_{n\ge0}J^n\). Par suite, on a:
\[
Y'=\Proj\bigl(i^*(S)\bigr)
=\Proj\left(\left(\bigoplus_{n\ge0}J^n\right)\otimes_{\cO_X}(\cO_X/J)\right)
=\Proj\left(\bigoplus_{n\ge0}J^n/J^{n+1}\right).
\]
Or, \(i\) étant régulière, \(\bigoplus_{n\ge0}J^n/J^{n+1}\simeq \underline{\Sym}_{\cO_Y}(J/J^2)\) (1.3). On note \(N\) le faisceau conormal \(J/J^2\); on a donc:
\[
Y'=\bfP(N).
\]

Sur \(Y'\), on peut écrire la suite exacte:
\begin{equation}\tag{3.1.2}
0\to F\to g^*(N)\to\cO_{Y'}(1)\to0.
\end{equation}

On remarque que \(\cO_{Y'}(1)=j^*(\cO_{X'}(1))=J'/J'^2\); c'est donc le faisceau conormal de l'immersion \(j\), et on le notera \(L\).

\textbf{Lemme 3.2.} Avec les hypothèses et les notations de 3.1, on a les isomorphismes suivants:
\[
\begin{array}{ll}
\text{i)}&
\forall i\ge0,\quad
\underline{\Tor}_i^{\cO_X}(\cO_{X'},\cO_Y)
\simeq
\bigwedge^i\underline{\Tor}_1^{\cO_X}(\cO_{X'},\cO_Y);\\[4pt]
\text{ii)}&
\underline{\Tor}_1^{\cO_X}(\cO_{X'},\cO_Y)\simeq F.
\end{array}
\]


